% GENERATED FILE. Edit canonical lessons, then run npm run generate:books.
% canonical-source-hash: fnv1a64:dfe1281f34580919
% canonical-lessons: BN-C466-shorir, BN-C466-bahu, BN-C466-angul, BN-C466-pith, BN-C466-ghar

\chapter{The Body}
\label{ch:bn-the-body}
\hlchaptermodality{466}

\begin{chapteropening}
\textbf{By the end of this chapter:} \emph{I can name parts of the body: the arms, fingers, back and neck.}

Complete the last lesson of chapter 466: I can name parts of the body: the arms, fingers, back and neck.
\end{chapteropening}

\section[\texorpdfstring{\bn{শরীর} (shôrīr)}{শরীর (shôrīr)}]{\bn{শরীর} (shôrīr) — the body}

\label{lesson:BN-C466-shorir}



\begin{quote}
Before the new one: say the Bengali for grateful, then the Bengali for to weigh up.
\end{quote}

\subsection*{You'll want to know: \bn{শরীর}}
\textbf{\bn{শরীর}} — \emph{shôrīr} — \textquotedblleft{}the body\textquotedblright{}.

\begin{quote}
\bn{সারাদিন} \bn{কাজ} \bn{করে} \bn{আমার} \bn{শরীর} \bn{ক্লান্ত।}
\end{quote}

\begin{quote}
\emph{sārādin kāj kôre āmār shôrīr klāntô.}
\end{quote}

\begin{quote}
After working all day my body is tired.
\end{quote}

\bn{শরীর} \bn{ভালো} \bn{নেই}, I'm not feeling well; \bn{শরীর} \bn{খারাপ}, unwell.

\begin{cousinweb}
\textbf{\bn{শরীর}} (\emph{shôrīr}) is a tatsama, Sanskrit \emph{śarīra}, \textquotedblleft{}the body\textquotedblright{} (Monier-Williams).
\end{cousinweb}

\begin{grammarlens}[title={What you've built}]
One more word for describing what you lived and felt.
\end{grammarlens}

\subsection*{Your turn}
\begin{itemize}
\raggedright
  \item \emph{Say it:} \emph{shôrīr}
  \item \emph{Say it:} \emph{shôrīr}, once more
  \item \emph{Say it:} \emph{sārādin kāj kôre āmār shôrīr klāntô.} — After working all day my body is tired.
  \item \emph{Recall:} say the Bengali for grateful, then the Bengali for to weigh up, then say \emph{shôrīr} again
\end{itemize}

\begin{tcolorbox}[breakable,colback=teal!4,colframe=teal!35!black,title={Before you move on}]
What does \bn{শরীর} mean? (\textquotedblleft{}The body\textquotedblright{}.) Say it once more.
\end{tcolorbox}
\section[\texorpdfstring{\bn{বাহু} (bāhu)}{বাহু (bāhu)}]{\bn{বাহু} (bāhu) — an arm (from shoulder to hand)}

\label{lesson:BN-C466-bahu}



\begin{quote}
Before the new one: say the Bengali for to weigh up, then the Bengali for the body.
\end{quote}

\subsection*{You'll want to know: \bn{বাহু}}
\textbf{\bn{বাহু}} — \emph{bāhu} — \textquotedblleft{}an arm (from shoulder to hand)\textquotedblright{}.

\begin{quote}
\bn{মীরা} \bn{শিশুটিকে} \bn{দুই} \bn{বাহুতে} \bn{তুলে} \bn{নিল।}
\end{quote}

\begin{quote}
\emph{mīrā shishuṭike dui bāhute tule nilô.}
\end{quote}

\begin{quote}
Mira lifted the child up in both arms.
\end{quote}

\bn{হাত} (taught) is the hand and, loosely, the whole arm; \bn{বাহু} is the arm itself.

\begin{cousinweb}
\textbf{\bn{বাহু}} (\emph{bāhu}) is a tatsama, Sanskrit \emph{bāhu}, \textquotedblleft{}arm\textquotedblright{} (CDIAL 9229), beside which CDIAL lists a worn-down Bengali \emph{bāh}.
\end{cousinweb}

\begin{grammarlens}[title={What you've built}]
One more word for describing what you lived and felt.
\end{grammarlens}

\subsection*{Your turn}
\begin{itemize}
\raggedright
  \item \emph{Say it:} \emph{bāhu}
  \item \emph{Say it:} \emph{bāhu}, once more
  \item \emph{Say it:} \emph{mīrā shishuṭike dui bāhute tule nilô.} — Mira lifted the child up in both arms.
  \item \emph{Recall:} say the Bengali for to weigh up, then the Bengali for the body, then say \emph{bāhu} again
\end{itemize}

\begin{tcolorbox}[breakable,colback=teal!4,colframe=teal!35!black,title={Before you move on}]
What does \bn{বাহু} mean? (\textquotedblleft{}An arm (from shoulder to hand)\textquotedblright{}.) Say it once more.
\end{tcolorbox}
\section[\texorpdfstring{\bn{আঙুল} (āngul)}{আঙুল (āngul)}]{\bn{আঙুল} (āngul) — a finger; a toe}

\label{lesson:BN-C466-angul}



\begin{quote}
Before the new one: say the Bengali for the body, then the Bengali for an arm.
\end{quote}

\subsection*{You'll want to know: \bn{আঙুল}}
\textbf{\bn{আঙুল}} — \emph{āngul} — \textquotedblleft{}a finger; a toe\textquotedblright{}.

\begin{quote}
\bn{ছুরিতে} \bn{মীরার} \bn{আঙুল} \bn{কাটল।}
\end{quote}

\begin{quote}
\emph{chhurite mīrār āngul kāṭlô.}
\end{quote}

\begin{quote}
Mira cut her finger on the knife.
\end{quote}

\begin{cousinweb}
\textbf{\bn{আঙুল}} (\emph{āngul}) comes from Sanskrit \emph{aṅguli}, \textquotedblleft{}finger\textquotedblright{} (CDIAL 135), which CDIAL lists in this Bengali form.
\end{cousinweb}

\begin{grammarlens}[title={What you've built}]
One more word for describing what you lived and felt.
\end{grammarlens}

\subsection*{Your turn}
\begin{itemize}
\raggedright
  \item \emph{Say it:} \emph{āngul}
  \item \emph{Say it:} \emph{āngul}, once more
  \item \emph{Say it:} \emph{chhurite mīrār āngul kāṭlô.} — Mira cut her finger on the knife.
  \item \emph{Recall:} say the Bengali for the body, then the Bengali for an arm, then say \emph{āngul} again
\end{itemize}

\begin{tcolorbox}[breakable,colback=teal!4,colframe=teal!35!black,title={Before you move on}]
What does \bn{আঙুল} mean? (\textquotedblleft{}A finger; a toe\textquotedblright{}.) Say it once more.
\end{tcolorbox}
\section[\texorpdfstring{\bn{পিঠ} (piṭh)}{পিঠ (piṭh)}]{\bn{পিঠ} (piṭh) — the back (of the body)}

\label{lesson:BN-C466-pith}



\begin{quote}
Before the new one: say the Bengali for an arm, then the Bengali for a finger.
\end{quote}

\subsection*{You'll want to know: \bn{পিঠ}}
\textbf{\bn{পিঠ}} — \emph{piṭh} — \textquotedblleft{}the back (of the body)\textquotedblright{}.

\begin{quote}
\bn{ভারী} \bn{ব্যাগটা} \bn{তুলে} \bn{বাবার} \bn{পিঠে} \bn{ব্যথা} \bn{হল।}
\end{quote}

\begin{quote}
\emph{bhārī byāgṭā tule bābār piṭhe bæthā hôlô.}
\end{quote}

\begin{quote}
Lifting the heavy bag gave my father a backache.
\end{quote}

\begin{cousinweb}
\textbf{\bn{পিঠ}} (\emph{piṭh}) is listed by CDIAL for Bengali, as the back, under Sanskrit \emph{pṛṣṭi} (CDIAL 8370).
\end{cousinweb}

\begin{grammarlens}[title={What you've built}]
One more word for describing what you lived and felt.
\end{grammarlens}

\subsection*{Your turn}
\begin{itemize}
\raggedright
  \item \emph{Say it:} \emph{piṭh}
  \item \emph{Say it:} \emph{piṭh}, once more
  \item \emph{Say it:} \emph{bhārī byāgṭā tule bābār piṭhe bæthā hôlô.} — Lifting the heavy bag gave my father a backache.
  \item \emph{Recall:} say the Bengali for an arm, then the Bengali for a finger, then say \emph{piṭh} again
\end{itemize}

\begin{tcolorbox}[breakable,colback=teal!4,colframe=teal!35!black,title={Before you move on}]
What does \bn{পিঠ} mean? (\textquotedblleft{}The back (of the body)\textquotedblright{}.) Say it once more.
\end{tcolorbox}
\section[\texorpdfstring{\bn{ঘাড়} (ghāṛ)}{ঘাড় (ghāṛ)}]{\bn{ঘাড়} (ghāṛ) — the back of the neck; the nape}

\label{lesson:BN-C466-ghar}



\begin{quote}
Before the new one: say the Bengali for a finger, then the Bengali for the back.
\end{quote}

\subsection*{You'll want to know: \bn{ঘাড়}}
\textbf{\bn{ঘাড়}} — \emph{ghāṛ} — \textquotedblleft{}the back of the neck; the nape\textquotedblright{}.

\begin{quote}
\bn{সারাদিন} \bn{ল্যাপটপে} \bn{কাজ} \bn{করে} \bn{অরুণের} \bn{ঘাড়ে} \bn{ব্যথা।}
\end{quote}

\begin{quote}
\emph{sārādin læpṭôpe kāj kôre ôruner ghāṛe bæthā.}
\end{quote}

\begin{quote}
After a day's work on the laptop, Arun's neck aches.
\end{quote}

\begin{cousinweb}
\textbf{\bn{ঘাড়}} (\emph{ghāṛ}) is listed by CDIAL for Bengali under Sanskrit \emph{ghāṭā}, \textquotedblleft{}the back of the neck\textquotedblright{} (CDIAL 4459).
\end{cousinweb}

\begin{grammarlens}[title={What you've built}]
One more word for describing what you lived and felt.
\end{grammarlens}

\subsection*{Your turn}
\begin{itemize}
\raggedright
  \item \emph{Say it:} \emph{ghāṛ}
  \item \emph{Say it:} \emph{ghāṛ}, once more
  \item \emph{Say it:} \emph{sārādin læpṭôpe kāj kôre ôruner ghāṛe bæthā.} — After a day's work on the laptop, Arun's neck aches.
  \item \emph{Recall:} say the Bengali for a finger, then the Bengali for the back, then say \emph{ghāṛ} again
\end{itemize}

\begin{tcolorbox}[breakable,colback=teal!4,colframe=teal!35!black,title={Before you move on}]
What does \bn{ঘাড়} mean? (\textquotedblleft{}The back of the neck; the nape\textquotedblright{}.) Say it once more.
\end{tcolorbox}
